\section{المكونات العتادية لنظام التقاط الصوت}
يتطلب بناء نظام معالجة صوتية مدمج اختيار مكونات عتادية تتناسب مع متطلبات المعالجة في الزمن الحقيقي (\textenglish{Real-Time Processing}) لضمان التقاط الإشارة بأقل نسبة ضجيج ممكنة وتفسيرها لحظياً. يتألف العتاد الصلب للنظام من جزأين رئيسيين:

\subsection{وحدة التقاط الصوت}
تُستخدم وحدة التقاط صوتية (\textenglish{Microphone}) مخصصة وذات حساسية عالية لالتقاط الترددات الصوتية البشرية المعيارية. تقوم هذه الوحدة بتحويل التغيرات الميكانيكية في ضغط الهواء الناتجة عن الكلام إلى إشارات كهربائية تناظرية صغيرة. ولضمان جودة الإشارة، يتم تمريرها عبر مرشح ممرر حزمة (\textenglish{Band-pass Filter}) لإزالة الضجيج عالي التردد والترددات المنخفضة المصاحبة للمحيط، تليها مرحلة التحويل الرقمي التناظري (\textenglish{ADC}) للحصول على عينات رقمية بمعدل أخذ عينات (\textenglish{Sampling Rate}) قياسي يبلغ $16\text{ kHz}$، وهو التردد الأمثل لعمل خوارزميات التعرف على الكلام الحديثة.

\subsection{وحدة المعالجة المركزية}
تعتمد البنية على وحدة معالجة مدمجة ذات أداء حسابي عالٍ وقدرة على استضافة النماذج اللغوية والصوتية للعمل بصورة مستقلة. تتطلب هذه الوحدة توافر ذاكرة وصول عشوائي (\textenglish{RAM}) كافية ومسرعات حسابية قادرة على تنفيذ العمليات المصفوفية المعقدة لشبكات العصبونية الاصطناعية المستخدمة في استخلاص الميزات الصوتية، وذلك لضمان زمن تأخير (\textenglish{Latency}) أدنى ما يمكن بين نطق الأمر وبين توليد الاستجابة الحركية الموافقة في مفاصل الذراع.